% TOP_PROBLEM: {"id": 3396, "title": "Exponential Algorithms for Knapsack", "category_id": 15, "category": {"id": 15, "name": "computer_science", "display_name": "Computer Science", "description": "Computational complexity, algorithms, and theoretical CS.", "slug": "computer-science", "order_index": 15, "created_at": "2026-07-31T15:26:25.671Z"}, "status": "open", "problem_number": "OPG-36311", "set_id": 12}
% FIRST_POSTED: 2026-10-01
% ENTRY_KIND: statement_only
% UnsolvedMath ID 3396; source catalogue code OPG-36311
% !TeX program = lualatex
% Definition and sources only; this file is not an LLM solution attempt.
\documentclass[11pt,a4paper]{article}
\usepackage[margin=25mm]{geometry}
\usepackage{fontspec}
\setmainfont{lmroman10-regular.otf}[BoldFont=lmroman10-bold.otf,
 ItalicFont=lmroman10-italic.otf,BoldItalicFont=lmroman10-bolditalic.otf]
\usepackage{amsmath,amssymb,mathtools}
\usepackage{unicode-math}
\setmathfont{latinmodern-math.otf}
\AtBeginDocument{\let\setminus\smallsetminus}
\usepackage{xurl,hyperref}
\hypersetup{colorlinks=true,urlcolor=blue}
\setcounter{secnumdepth}{0}
\setlength{\parindent}{0pt}
\setlength{\parskip}{5pt}
\setlength{\emergencystretch}{3em}
\begin{document}
\section*{Exponential Algorithms for Knapsack}\label{3396}
{\small UnsolvedMath ID 3396 · OPG-36311 · Computer Science · OpenGarden}
\subsection{Definitions and mathematical statement}
The input is a list of integers $a_1,\ldots,a_n,b$, encoded in binary.
The decision problem specified by the source is subset sum:
\[
 \exists (x_1,\ldots,x_n)\in\{0,1\}^n\quad
                        \sum_{i=1}^n a_ix_i=b.
\]
Each list position can be used at most once; equal integer values
may appear in different positions. The empty selection is allowed.
Although the source calls this knapsack, there is no additional
profit objective or inequality capacity constraint.

Let $L$ be the total bit length of the input, including signs. The
algorithmic target is a uniform exact worst-case algorithm using
\[
                            O(2^{n/3}L^c)
\]
bit operations for some absolute constant $c$, on every input.
This states the source's exponential rate while making the omitted
cost of reading and handling large integers explicit. No space bound
is imposed. Deterministic algorithms are the baseline interpretation;
any randomized variant must state its error and time guarantees
separately. A faster method that works only for random or specially
structured integer lists does not meet the universal worst-case
target. The problem is to achieve this running time or to establish
a rigorously specified obstruction in an explicit computational model.
\subsection{Short English statement}
Can arbitrary binary-encoded subset-sum instances on n integers be solved exactly in time two to the n/3 times a polynomial in the input length?
\subsection{Sources}
\begin{itemize}
\item[S1] ULAM.AI and the UnsolvedMath Contributors, supplied problems.json, record 3396 (OPG-36311), OpenGarden collection. Catalogue identity and source transcription; CC BY 4.0.
\url{https://huggingface.co/datasets/ulamai/UnsolvedMath}.
\item[S2] Open Problem Garden, Exponential Algorithms for Knapsack. Original problem and discussion; the mathematical formulation below is expanded from the supplied source record.
\url{http://www.openproblemgarden.org/op/exponential_algorithms_for_knapsack}.
\end{itemize}
\end{document}
